% Source: https://artofproblemsolving.com/community/c3h1787556_pls_help
% Stable identity: aops-topic-1787556-q01
\textbf{$A$.} With $y=\dfrac{1}{3}$: $3x+3=2$, so $x=-\dfrac{1}{3}$ and $A=\left(-\dfrac{1}{3},\dfrac{1}{3}\right)$.

\textbf{$B$.} With $x=\dfrac{2}{3}$: $2+9y=2$, so $y=0$ and $B=\left(\dfrac{2}{3},0\right)$.

\textbf{Slope.} Through $A$ and $B$,
\[
m=\dfrac{0-\dfrac{1}{3}}{\dfrac{2}{3}-\left(-\dfrac{1}{3}\right)}=\dfrac{-\dfrac{1}{3}}{1}=\boxed{-\dfrac{1}{3}}.
\]
(Equivalently, $y=-\dfrac{1}{3}x+\dfrac{2}{9}$.)

\begin{asy}
usepackage("amsmath");
size(7cm);
defaultpen(fontsize(10pt));
import graph;
real f(real x){return (2-3x)/9;}
draw((-1.7,0)--(1.4,0),Arrow(6));draw((0,-0.45)--(0,0.8),Arrow(6));
label("$x$",(1.4,0),E);label("$y$",(0,0.8),N);
draw(graph(f,-1.6,1.3),blue+linewidth(0.9));
pair A=(-1/3,1/3),B=(2/3,0);
dot(A);dot(B);
draw((-1/3,0)--A--(0,1/3),dashed+gray);
label("$A$",A,2*dir(215));label("$-\dfrac{1}{3}$",(-1/3,0),S);draw((0,1/3)--(0.04,1/3));label("$\dfrac{1}{3}$",(0.03,1/3),dir(20),fontsize(9pt));
label("$B$",B,SW);label("$\dfrac{2}{3}$",(2/3,0),2*N);
label("$3x+9y=2$",(1.3,f(1.3)),SE,blue);
label("\mbox{}",(min(currentpicture,user=true)+max(currentpicture,user=true))/2,black); // keeps the text after the figure black when LaTeX typesets the labels inline
\end{asy}
